% GENERATED FILE. Edit canonical lessons, then run npm run generate:books.
% canonical-source-hash: fnv1a64:2e7b917839bcdbf3
% canonical-lessons: LA-C408-mendum, LA-C408-vitium, LA-C408-ferula, LA-C408-verber, LA-C408-vapulo

\chapter{Mistakes and the Cane}
\label{ch:la-mistakes-and-the-cane}
\hlchaptermodality{408}

\begin{chapteropening}
\textbf{By the end of this chapter:} \emph{I can talk about a mistake in a text and a fault, the cane and the lash, and say that a pupil got a beating.}

Complete the last lesson of chapter 408: I can talk about a mistake in a text and a fault, the cane and the lash, and say that a pupil got a beating.
\end{chapteropening}

\section[mendum, mendī]{mendum, mendī — a mistake, a blemish (in writing)}

\label{lesson:LA-C408-mendum}



\begin{quote}
Before the new one: say the Latin for ignorance, then the Latin for to imitate.
\end{quote}

\subsection*{You'll want to know: mendum, mendī}
\textbf{mendum, mendī} — \textquotedblleft{}a mistake, a blemish (in writing)\textquotedblright{}.

\begin{quote}
In hōc librō multa menda sunt.
\end{quote}

\begin{quote}
\emph{There are many mistakes in this book.}
\end{quote}

\emph{mendum, mendī} n. \emph{Ēmendō}, to correct, is to take the mistakes out. English \emph{amend} and \emph{emend}.

\begin{cousinweb}
\emph{mendum} is inherited from a Proto-Indo-European root meaning \textquotedblleft{}a fault, a flaw\textquotedblright{} (EtymDB), and \emph{ēmendō} \textquotedblleft{}to correct\textquotedblright{} is built on it (Lewis \& Short). English \emph{amend} and \emph{emend} come from \emph{ēmendō}.
\end{cousinweb}

\begin{grammarlens}[title={What you've built}]
One more word for telling, explaining and describing at B1.
\end{grammarlens}

\subsection*{Your turn}
\begin{itemize}
\raggedright
  \item \emph{Say it:} \emph{mendum, mendī}
  \item \emph{Say it:} \emph{mendum, mendī}, once more
  \item \emph{Say it:} a sentence of your own with \emph{mendum, mendī}
  \item \emph{Recall:} say the Latin for ignorance, then the Latin for to imitate, then say \emph{mendum, mendī} again
\end{itemize}

\begin{tcolorbox}[breakable,colback=teal!4,colframe=teal!35!black,title={Before you move on}]
What does mendum, mendī mean? (\textquotedblleft{}A mistake, a blemish (in writing)\textquotedblright{}.) Say it once more.
\end{tcolorbox}
\section[vitium, vitiī]{vitium, vitiī — a fault, a defect; a vice}

\label{lesson:LA-C408-vitium}



\begin{quote}
Before the new one: say the Latin for to imitate, then the Latin for a mistake.
\end{quote}

\subsection*{You'll want to know: vitium, vitiī}
\textbf{vitium, vitiī} — \textquotedblleft{}a fault, a defect; a vice\textquotedblright{}.

\begin{quote}
Nēmō sine vitiīs nāscitur.
\end{quote}

\begin{quote}
\emph{No one is born without faults.}
\end{quote}

\emph{vitium, vitiī} n. A flaw in a thing, a fault in a person. English \emph{vice} and \emph{vicious}.

\begin{cousinweb}
The origin of \emph{vitium} is not settled: Lewis \& Short tie it to the root of \emph{vieō} \textquotedblleft{}to twist\textquotedblright{}, so first a twist, a flaw, while EtymDB offers a different Proto-Indo-European source. English \emph{vice} and \emph{vicious} come from it.
\end{cousinweb}

\begin{grammarlens}[title={What you've built}]
One more word for telling, explaining and describing at B1.
\end{grammarlens}

\subsection*{Your turn}
\begin{itemize}
\raggedright
  \item \emph{Say it:} \emph{vitium, vitiī}
  \item \emph{Say it:} \emph{vitium, vitiī}, once more
  \item \emph{Say it:} a sentence of your own with \emph{vitium, vitiī}
  \item \emph{Recall:} say the Latin for to imitate, then the Latin for a mistake, then say \emph{vitium, vitiī} again
\end{itemize}

\begin{tcolorbox}[breakable,colback=teal!4,colframe=teal!35!black,title={Before you move on}]
What does vitium, vitiī mean? (\textquotedblleft{}A fault, a defect; a vice\textquotedblright{}.) Say it once more.
\end{tcolorbox}
\section[ferula, ferulae]{ferula, ferulae — a cane (for punishing pupils); a fennel stalk}

\label{lesson:LA-C408-ferula}



\begin{quote}
Before the new one: say the Latin for a mistake, then the Latin for a fault.
\end{quote}

\subsection*{You'll want to know: ferula, ferulae}
\textbf{ferula, ferulae} — \textquotedblleft{}a cane (for punishing pupils); a fennel stalk\textquotedblright{}.

\begin{quote}
Magister ferulam in manū tenēbat.
\end{quote}

\begin{quote}
\emph{The teacher held the cane in his hand.}
\end{quote}

\emph{ferula, ferulae} f.: the stalk of the giant fennel, light and stinging, the schoolmaster's cane. English \emph{ferule}.

\begin{cousinweb}
The origin of \emph{ferula} is not settled: EtymDB derives it from \emph{feriō} \textquotedblleft{}to strike\textquotedblright{}, but Lewis \& Short give no origin for it. English \emph{ferule} comes from it.
\end{cousinweb}

\begin{grammarlens}[title={What you've built}]
One more word for telling, explaining and describing at B1.
\end{grammarlens}

\subsection*{Your turn}
\begin{itemize}
\raggedright
  \item \emph{Say it:} \emph{ferula, ferulae}
  \item \emph{Say it:} \emph{ferula, ferulae}, once more
  \item \emph{Say it:} a sentence of your own with \emph{ferula, ferulae}
  \item \emph{Recall:} say the Latin for a mistake, then the Latin for a fault, then say \emph{ferula, ferulae} again
\end{itemize}

\begin{tcolorbox}[breakable,colback=teal!4,colframe=teal!35!black,title={Before you move on}]
What does ferula, ferulae mean? (\textquotedblleft{}A cane (for punishing pupils); a fennel stalk\textquotedblright{}.) Say it once more.
\end{tcolorbox}
\section[verber, verberis]{verber, verberis — a lash, a blow; (plural) a beating}

\label{lesson:LA-C408-verber}



\begin{quote}
Before the new one: say the Latin for a fault, then the Latin for a cane.
\end{quote}

\subsection*{You'll want to know: verber, verberis}
\textbf{verber, verberis} — \textquotedblleft{}a lash, a blow; (plural) a beating\textquotedblright{}.

\begin{quote}
Puer verbera timet.
\end{quote}

\begin{quote}
\emph{The boy fears a beating.}
\end{quote}

\emph{verber, verberis} n., mostly plural: \emph{verbera}, blows, a flogging. English \emph{reverberate} comes from its verb.

\begin{cousinweb}
\emph{verber} is inherited from a Proto-Indo-European root meaning \textquotedblleft{}to turn, to bend\textquotedblright{} (EtymDB), which also gives English \emph{warp}. English \emph{reverberate} comes from its verb \emph{verberō} \textquotedblleft{}to beat\textquotedblright{}.
\end{cousinweb}

\begin{grammarlens}[title={What you've built}]
One more word for telling, explaining and describing at B1.
\end{grammarlens}

\subsection*{Your turn}
\begin{itemize}
\raggedright
  \item \emph{Say it:} \emph{verber, verberis}
  \item \emph{Say it:} \emph{verber, verberis}, once more
  \item \emph{Say it:} a sentence of your own with \emph{verber, verberis}
  \item \emph{Recall:} say the Latin for a fault, then the Latin for a cane, then say \emph{verber, verberis} again
\end{itemize}

\begin{tcolorbox}[breakable,colback=teal!4,colframe=teal!35!black,title={Before you move on}]
What does verber, verberis mean? (\textquotedblleft{}A lash, a blow; (plural) a beating\textquotedblright{}.) Say it once more.
\end{tcolorbox}
\section[vāpulō, vāpulāre]{vāpulō, vāpulāre — to be beaten, to get a thrashing}

\label{lesson:LA-C408-vapulo}



\begin{quote}
Before the new one: say the Latin for a cane, then the Latin for a lash.
\end{quote}

\subsection*{You'll want to know: vāpulō, vāpulāre}
\textbf{vāpulō, vāpulāre} — \textquotedblleft{}to be beaten, to get a thrashing\textquotedblright{}.

\begin{quote}
Puer piger vāpulāvit.
\end{quote}

\begin{quote}
\emph{The lazy boy got a beating.}
\end{quote}

Active in form but passive in meaning: \emph{vāpulō} means 'I am beaten'. Its perfect is \emph{vāpulāvī}. \emph{Vāpulābis!}, you'll get a thrashing, is a threat in comedy.

\begin{cousinweb}
The origin of \emph{vāpulō} is not settled: Lewis \& Short suggest, with a \textquotedblleft{}perhaps\textquotedblright{}, a root meaning to wriggle, and EtymDB gives a Proto-Indo-European root without a meaning.
\end{cousinweb}

\begin{grammarlens}[title={What you've built}]
One more word for telling, explaining and describing at B1.
\end{grammarlens}

\subsection*{Your turn}
\begin{itemize}
\raggedright
  \item \emph{Say it:} \emph{vāpulō, vāpulāre}
  \item \emph{Say it:} \emph{vāpulō, vāpulāre}, once more
  \item \emph{Say it:} a sentence of your own with \emph{vāpulō, vāpulāre}
  \item \emph{Recall:} say the Latin for a cane, then the Latin for a lash, then say \emph{vāpulō, vāpulāre} again
\end{itemize}

\begin{tcolorbox}[breakable,colback=teal!4,colframe=teal!35!black,title={Before you move on}]
What does vāpulō, vāpulāre mean? (\textquotedblleft{}To be beaten, to get a thrashing\textquotedblright{}.) Say it once more.
\end{tcolorbox}
